%! TeX root = ../main.tex
\chapter{TINJAUAN PUSTAKA}
\label{chap:tinjauanpustaka}

\section{Hasil Penelitian Terdahulu}
\label{sec:penelitianterdahulu}

Pada bagian ini dijelaskan kajian dari penelitian-penelitian terdahulu yang relevan dengan topik tugas akhir. Perkembangan teknologi kecerdasan buatan dalam bidang kesehatan telah berkembang secara pesat dalam mendukung efisiensi pelayanan klinis \parencite{rajpurkar2022ai}.

\lipsum[1]

\section{Dasar Teori}
\label{sec:dasarteori}

Bagian ini memuat teori-teori dasar, konsep, dan algoritma yang melandasi penelitian.

\subsection{Teori Pendukung Pertama}
\label{subsec:teoripertama}

\lipsum[2]

\begin{equation}
  \label{eq:persamaan1}
  \sum \mathbf{F} = 0\; \Leftrightarrow\; \frac{\mathrm{d} \mathbf{v} }{\mathrm{d}t} = 0
\end{equation}

\subsection{Teori Pendukung Kedua}
\label{subsec:teorikedua}

\lipsum[3]
